% ==========================================================
% STANDARDIZED PREAMBLE (STATIC CORE) — DO NOT EDIT THIS PART
% ==========================================================

\documentclass[11pt,a4paper]{article}

\usepackage{iftex}
\ifPDFTeX
  \usepackage[utf8]{inputenc}
  \usepackage[T1]{fontenc}
  \usepackage{stix2}
\else
  \usepackage{fontspec}
  \usepackage{unicode-math}
  \setmainfont{STIX Two Text}
  \setmathfont{STIX Two Math}
\fi

\usepackage[margin=2.5cm]{geometry}
\usepackage{parskip}
\usepackage{microtype}
\usepackage{enumitem}

\usepackage{amsmath, amssymb}

\usepackage{tocloft}
\setlength{\cftbeforesecskip}{2pt}
\setlength{\cftbeforesubsecskip}{1pt}

\widowpenalty=10000
\clubpenalty=10000

\usepackage{xcolor}
\definecolor{myblue}{RGB}{0, 51, 153}

\usepackage{sectsty}
\partfont{\color{myblue}}
\chapterfont{\color{myblue}}
\sectionfont{\color{myblue}}
\subsectionfont{\color{myblue}}
\subsubsectionfont{\color{myblue}}

\usepackage{hyperref}

\renewcommand{\cftsecfont}{\bfseries\color{myblue}}
\renewcommand{\cftsubsecfont}{\color{myblue}}
\renewcommand{\cftsubsubsecfont}{\color{myblue}}
\renewcommand{\cftsecpagefont}{\bfseries}
\renewcommand{\cftsubsecpagefont}{}
\renewcommand{\cftsubsubsecpagefont}{}

\hypersetup{
  colorlinks=true,
  linkcolor=myblue,
  citecolor=myblue,
  urlcolor=myblue
}

\usepackage{titlesec}
\titlespacing{\section}{0pt}{12pt}{6pt}
\titlespacing{\subsection}{0pt}{10pt}{5pt}
\titlespacing*{\section}{0pt}{2.5ex plus 1ex minus .2ex}{1.5ex plus .2ex}
\titlespacing*{\subsection}{0pt}{2ex plus 0.8ex minus .2ex}{1.2ex plus .2ex}
\titlespacing*{\subsubsection}{0pt}{1.8ex plus 0.6ex minus .2ex}{1ex plus .2ex}

\makeatletter
\newcommand{\StartBlueDoc}[3][]{%
  \title{\vspace{-2em}\textbf{#2}%
    \if\relax\detokenize{#3}\relax\else\\[2pt]{\large #3}\fi
    \vspace{+0.6em}}
  \author{#1}
  \date{\vspace{-0.6em}\today}
  \begin{document}
  \maketitle
  \tableofcontents
  \vspace{1em}\hrule\vspace{1em}
}
\makeatother

% ==========================================================
% CUSTOMIZATION ZONE
% ==========================================================

\newcommand{\DocAuthor}{}
\newcommand{\DocTitle}{Design Principles for Scientific Software}
\newcommand{\DocSubtitle}{}

\setcounter{tocdepth}{2}

\usepackage{tikz}

% ==================
% Start the document
% ==================
\StartBlueDoc[\DocAuthor]{\DocTitle}{\DocSubtitle}



% ==========================================================
\section*{About these guidelines}
\phantomsection
\addcontentsline{toc}{section}{About these guidelines}
% ==========================================================

These guidelines apply wherever science is turned into software that several
people use over and over, whether it is published openly or kept within an
organization, and independent of programming language and scientific field.
Usually the science is known; the task is to make it reliable, automated, and
fit to run at scale. One-off analysis scripts are out of scope.



\subsection*{Normative language}
\phantomsection
\addcontentsline{toc}{subsection}{Normative language}
\label{sec:normative}

The guidelines distinguish requirements from defaults. The following words
carry a fixed meaning throughout, in the spirit of RFC~2119:

\begin{description}
  \item[must, must not] A requirement. Code that violates it is a defect and
    is fixed, not justified.
  \item[should, should not] A strong default. A deviation is acceptable only
    for a concrete reason, which is stated where the deviation occurs, for
    example in a code comment or in the review. The role of the software
    (Subsection~\ref{sec:role}) is such a reason; it is stated once for the
    whole package, for example in its documentation, rather than at each
    deviation.
  \item[may] Permitted, but not required.
\end{description}

Plain imperatives such as ``prefer'' and ``avoid'' carry the weight of
\emph{should}. When two rules pull in different directions, the priorities in
Subsection~\ref{sec:quality} decide.



\subsection*{Terminology}
\phantomsection
\addcontentsline{toc}{subsection}{Terminology}
\label{sec:terminology}

\begin{description}
  \item[Package] The software these guidelines apply to, whatever its form: a
    library, a command-line tool, or a workflow.
  \item[Library code] The part of a package that other code calls: the
    scientific logic and everything it needs. Most rules in these guidelines
    are about library code.
  \item[Application layer] The part of a tool or workflow that owns the
    program: command-line entry points, workflow definitions, and job
    scripts. Rules that restrict library code, such as those on printing,
    files, and global state, do not apply to it; instead, it has to stay
    thin (Subsection~\ref{sec:library}).
  \item[Host stack] The software ecosystem the user already works in: the
    programming language, its standard array and data-frame types, its
    established numerical, statistical, and plotting libraries, and the
    user's own workflow tools. Examples are the scientific Python ecosystem,
    R with its package repositories, and Julia with its package ecosystem.
    The host stack sets the conventions for data types, randomness,
    warnings, logging, and packaging that the package follows.
  \item[User-facing function] Any function, class, or method that is
    documented for users, including low-level primitives. Together, the
    user-facing functions form the public API; everything else is internal.
  \item[Caller] The code that calls a user-facing function, and by extension
    the user who wrote it.
\end{description}



% ==========================================================
\section{Guiding priorities}
% ==========================================================

\subsection{Build the right thing before building the thing right}

Establish that the overall approach is correct before investing in
optimization, testing, documentation, or polish. Effort spent in the wrong
direction is not partially recovered; it is lost in full, and the better the
execution, the greater the loss.

Incremental improvement converges on the nearest local optimum, not on the best
achievable one. A design can therefore become steadily better while remaining
fundamentally misplaced, and no amount of further refinement will reveal this,
because refinement measures progress against the current approach rather than
against the alternatives to it.

Quality and speed multiply direction rather than substitute for it. Executing a
wrong approach quickly and well only arrives sooner at a place that must be
abandoned. Refinement also entrenches: the more effort a design has absorbed,
the harder it becomes to question, so the assessment should happen early and be
repeated deliberately.

Identify which decisions are reversible and which are not. Reversible decisions
can be made quickly and corrected later. Irreversible decisions, such as the
core abstraction and the public interface, deserve the effort that would
otherwise be spent perfecting details, because they determine whether that
effort is worth anything at all.



\subsection{Get your priorities right}
\label{sec:quality}

When goals conflict, rank them as follows, and let the higher one win:

\begin{enumerate}
  \item \textbf{Correctness.} The results are right and reproducible; a
    result that cannot be reproduced cannot be checked
    (Subsection~\ref{sec:reproducibility}).
  \item \textbf{Clarity.} Code, interface, and behavior can be understood.
    Code that cannot be understood hides its errors.
  \item \textbf{Testability.} Correctness can be demonstrated, not merely
    claimed.
  \item \textbf{Maintainability.} The package stays correct and clear as it
    changes.
  \item \textbf{Performance.} The package is fast enough for the data sizes
    of real applications (Subsection~\ref{sec:size-matters}).
  \item \textbf{Features.} More functionality, options, and code are a
    bonus, never a substitute for the priorities above.
\end{enumerate}

The same ranking resolves conflicts between the rules in these guidelines:
when two rules pull in different directions, the rule that serves the higher
priority wins. Prefer a small amount of reliable functionality over a large
amount of unreliable functionality; code that is produced quickly but is hard
to understand, test, or trust only adds to the burden of the package.



\subsection{Better loud than wrong}
\label{sec:loud}

When something goes wrong, the package must fail visibly rather than return a
result that looks valid but cannot be trusted. Given the choice between an
error and a result that may be wrong, choose the error.

The two outcomes are not symmetric. A failure is annoying and can cost time,
but it is visible and its cost is bounded: the user sees it, fixes the cause,
and runs again. A silently wrong result looks exactly like a correct one. It
flows into further analyses, figures, publications, and decisions, and nobody
looks for the error, because nothing suggests that there is one. A failure
costs convenience; a silently wrong result costs correctness, which ranks
first (Subsection~\ref{sec:quality}).

Silent failures are rarely designed. They usually arise from code that tries
to be helpful or robust. Typical sources are:

\begin{itemize}
  \item Catching an error and returning a default value, an empty result, or
    NaN in place of the result.
  \item Falling back to a different algorithm, approximation, or data source
    when the intended one fails or an optional dependency is missing.
  \item Letting integers overflow, clipping or truncating values to fit a type
    or range, or broadcasting arrays whose shapes do not match.
  \item Dropping records that find no match in a join or alignment.
  \item Continuing after a step has failed, so that an incomplete output looks
    complete.
\end{itemize}

Each of these turns a failure the user would have seen into a wrong result
the user will not see. Let errors that cannot be handled correctly propagate,
and make every fallback, coercion, and omission explicit.

Being loud does not mean giving up at the first problem. A workflow that
processes thousands of samples may skip one that fails, and an iterative
method may return a result that did not converge
(Subsection~\ref{sec:convergence}), provided the degradation is explicit: it
is documented or chosen by the caller, the result records what is missing or
uncertain, and the failure is reported through the channels of
Subsection~\ref{sec:output}. A tool or workflow that did not complete what was
asked must also say so in its exit status. What must never happen is a
degraded result that cannot be told apart from a valid one.

Many rules in these guidelines apply this principle to a particular situation.
Where none of them covers a case, the default is the same: when in doubt,
fail.



\subsection{Pay as you grow}
\label{sec:role}

Engineering effort has to pay for itself. Not all software in scope deserves
the same investment: a tool that three colleagues run on small datasets needs
less than a workflow that processes every sample of a core facility, or a
library whose results enter publications and regulatory decisions. Scale the
effort with the role of the software, judged by how many people use it, how
often and for how long, on how much data, and what a wrong result would cost.

The requirements marked \emph{must} are the floor for all software in scope;
most of them cost little and prevent wrong results. What scales with the role
is everything above that floor: how thoroughly the \emph{should}s are
followed, the depth of testing and documentation, the effort put into
performance, and the length of the stability commitment.

Roles change, and software usually grows into a larger role without anyone
deciding it. Reassess the role whenever the users, the data, or the stakes
change, and raise the effort before new users rely on the software, not after
the first wrong result.


\subsection{KISS: Keep it simple, stupid}

Prefer the simplest design that solves the problem well. Simplicity is
measured by how much a reader has to hold in mind to understand the code and
its behavior: the number of concepts, options, code paths, and special cases,
not the number of lines.

Every option multiplies the behaviors that have to be documented and tested,
because options interact. Every layer of indirection is a place where a reader
has to stop and look elsewhere. Add either only when a concrete application
demands it (Subsection~\ref{sec:applications}).

Prefer plain, direct code over clever code. When two designs solve the problem
equally well, choose the one with fewer concepts. Simple does not mean
simplistic: a design that looks simple because it ignores part of the problem,
such as missing values or edge cases, is not simple but incomplete.



% ==========================================================
\section{Scope}
% ==========================================================

\subsection{Own only the core domain}
\label{sec:core-domain}

Custom code belongs only in the distinctive scientific problem that justifies
the package. That problem is the core domain. Linear algebra, optimization,
numerical integration, automatic differentiation, parallel execution, file
input and output, and plotting are generic capabilities: reuse mature
libraries from the host stack instead of rebuilding them.

A package must document, test, and maintain everything it owns. Owning generic
work therefore takes effort away from the scientific contribution, and a
reimplementation is usually less tested, less optimized, and less numerically
robust than the established one. Implement the core; compose with the host
stack for the rest.



\subsection{Code is for life, not just for Christmas}
\label{sec:sustainability}

Everything the package takes on is a commitment for the lifetime of the
software, and that lifetime usually exceeds the involvement of the people who
wrote it. Removing something later breaks the users who rely on it
(Subsection~\ref{sec:stable-api}). Judge each addition by what it will cost
to maintain over years, not by the effort of adding it today, and prefer
established, well-understood technology over the newest option.

\textbf{Features.} Every feature adds to what the package must document,
test, and maintain (Subsection~\ref{sec:core-domain}) for as long as it
lives. Prefer a smaller package with well-tested core functionality over a
larger one with many unfinished or poorly supported features.

\textbf{Dependencies.} Each dependency is code the package did not write but
still has to install, pin, test against, license, and debug. Prefer the host
stack the user already has, and add a dependency only when it owns a job the
package must not reimplement and cannot leave to the caller. Choose
dependencies that will outlive your own project: widely used, maintained by
more than one person, backed by an institution or an active community, and
with a history of regular, compatible releases. Avoid dependencies that rest
on a single maintainer, thesis, or grant; where a fragile one is unavoidable,
isolate it behind a thin internal interface so that it can be replaced.
Dependencies needed only for optional functionality, such as plotting, should
be optional, so that the core installs without them.

\textbf{The package itself.} Build it to survive its authors. No part of the
code should be understood by only one person, and the knowledge needed to
maintain it, such as design decisions, conventions, and the release process,
belongs in the repository, not in people's heads.



\subsection{Don't call us, we'll call you}
\label{sec:library}

A framework follows the Hollywood principle: it owns the program and calls
user code at fixed hooks. A library does the opposite: the user owns the
program and calls the library. The scientific logic of a package should be
library code.

Library code should not take over the main program, the solver or
simulation loop, the data model, or the analysis workflow. Provide
composable primitives that users call from their own scripts, notebooks,
pipelines, and host-stack tools. The package should not become a framework
on top of the host stack.

Where the method needs user code, such as an objective function, a model, or
the right-hand side of a differential equation, accept it as an ordinary
function argument. Requiring users to subclass package classes or to
register code with a package-managed lifecycle turns the library into a
framework.

A package that is itself a tool or a workflow does own its program, through
its application layer (Subsection~\ref{sec:terminology}). Keep that layer
thin: it orchestrates, handles files and configuration, and calls the
library code, but contains no scientific logic. The science can then be
tested, reused, and scaled without running the whole workflow.



\subsection{Operate on data objects, not files}
\label{sec:no-files}

Library code should take data objects and return data objects; reading and
writing files belongs to the caller or to the application layer
(Subsection~\ref{sec:terminology}). This avoids assumptions about file
formats, directory layouts, and environments, and it prevents hidden side
effects.

Accept the data types of the host stack, including lazy, memory-mapped, and
chunked representations where the method permits it. Operating on data
objects does not mean loading everything into memory; materializing such data
behind the user's back is a surprise cost (Subsection~\ref{sec:cost}).

Keeping input and output out of library code creates an obligation: the
package's own objects, such as fitted models or result containers, should be
storable with the host stack's standard serialization mechanisms, or
convertible to standard types that are. If reading a domain-specific file
format is itself part of the core domain, isolate it in a separate module
that returns standard data objects.



\subsection{Avoid plotting in core functionality}
\label{sec:plotting}

Plotting is brittle, backend-dependent, and parameter-heavy, so core
functionality should not plot. If the package provides plotting, strictly
separate computation from visualization: functions that compute return
results, and separate functions plot them.

Plotting functions should take computed results as input and never recompute
them. Where the plotting library supports it, they should draw into a target
the caller can pass in, return the plot object instead of displaying it, and
leave global plotting styles untouched (Subsection~\ref{sec:global-state}).
The plotting library should be an optional dependency
(Subsection~\ref{sec:sustainability}).



% ==========================================================
\section{Code structure and abstraction}
% ==========================================================

\subsection{Derive abstractions from applications, not applications from abstractions}
\label{sec:applications}

Build for concrete applications that someone actually wants to run. An
abstraction invented first, in the hope that useful applications will follow
from it, is a guess about demand, and such guesses fail in ways that become
visible only once real work is attempted.

Generality that no application requested is speculative generality. It adds
options, layers, and special cases that must be documented, tested, and
maintained forever while solving no problem anyone has. The cost is paid
immediately, the benefit remains hypothetical.

Several genuinely different use cases are usually required before the shared
concept becomes visible; the rule of three sets the customary threshold at
three. A single use case invites a design fitted too tightly to it, and no use
case invites a design fitted to nothing.

Keep the driving applications in the repository as runnable examples or
integration tests. They document intent, they fail when the design drifts away
from practice, and they force every proposed feature to name the concrete task
that justifies it. Software is a practical artifact and can only be judged in
use.



\subsection{Design modules around clear responsibilities}
\label{sec:modules}

Organize code into modules with clear, coherent responsibilities. Each module
should own one part of the package, such as validation, computation,
preprocessing, or visualization, and its name should say which. Avoid
miscellaneous helper modules such as \texttt{utils} when a more specific name
would describe the purpose of the code; a module that cannot be named
precisely usually has no single responsibility.

Dependencies between modules should point in one direction, from the specific
to the fundamental. Computation must not depend on visualization, and core
computation should not depend on input and output. Knowledge shared by several
modules, such as validation checks and numerical defaults, belongs in one
module that the others import (Subsection~\ref{sec:dry}).



\subsection{No code at large}
\label{sec:functions}

Library code lives in functions and classes. At module level, only imports,
named constants (Subsection~\ref{sec:magic-numbers}), and definitions belong.
A statement that does work there runs on every import
(Subsection~\ref{sec:import}), cannot be tested in isolation, and tends to
communicate through global variables. The application layer follows the same
rule: its logic goes into a main function, and the entry point only calls it.

Prefer functions. Most scientific operations take data and parameters and
return a result, which a function expresses directly and which is easy to test
and compose. Use a class when there is state with invariants to protect, such
as a fitted model, or a result that carries its values together with
diagnostics and provenance (Subsection~\ref{sec:traceability}). A class that
only groups functions, or stores its arguments for its methods to read back,
should be a function.



\subsection{DRY: Don't repeat yourself}
\label{sec:dry}

Every piece of knowledge should have a single, authoritative
representation in the software. DRY is not primarily about code.
Knowledge includes rules, assumptions, constants, numerical defaults,
conventions such as axes and units, validation logic, data schemas,
and the documentation that describes them.

When the same knowledge lives in several places, every change has to find
and update all copies consistently. Copies that drift apart produce
software that contradicts itself. Typical examples in scientific code
are a tolerance hard-coded in several functions, the same input check
reimplemented in each user-facing function, or a default value repeated in
the function signature, its documentation, and the examples. Define such
knowledge once and refer to it everywhere else. Where documentation has to
state a value, generate it from the single definition if the tooling permits.

DRY concerns knowledge, not text. Code that looks different can still
duplicate the same rule. Conversely, code that looks alike may encode
different rules, and merging it is a mistake
(see Subsection~\ref{sec:wrong-abstraction}).



\subsection{No magic numbers}
\label{sec:magic-numbers}

Every numerical constant with a scientific or numerical role, such as a
threshold, cutoff, tolerance, iteration limit, or physical constant, must have
a name, a unit where it has one (Subsection~\ref{sec:conventions}), a single
definition (Subsection~\ref{sec:dry}), and a justification stated next to that
definition: a reference, a derivation, or the reason for the choice. In
scientific code, an unexplained \texttt{3.0} or \texttt{0.05} is often an
unexamined modeling decision, and a reviewer cannot tell a deliberate choice
from an arbitrary one.

If a constant affects the results and users might reasonably want a different
value, expose it as a parameter; its default then becomes part of the API
(Subsection~\ref{sec:numeric-defaults}). Take physical constants from an
authoritative source in the host stack rather than typing them in by hand.
Constants whose meaning is evident from the code, such as the 2 in $2\pi$,
need no name.



\subsection{Duplicate code is cheaper than the wrong abstraction}
\label{sec:wrong-abstraction}

Avoid premature abstraction. Similar-looking code should be unified only when
it represents the same underlying concept or rule (Subsection~\ref{sec:dry}).
A wrong abstraction hides important differences, increases complexity, and
makes future changes harder, because every change has to respect all of its
callers at once. Prefer simple duplication until the shared idea is clear and
stable.

A common sign of a wrong abstraction is a shared function that accumulates
flags and parameters whose only purpose is to switch its behavior for
particular callers. When that happens, inline the function back into its
callers, remove from each copy what that caller does not need, and extract a
new abstraction only once the genuinely shared rule has become visible.



% ==========================================================
\section{API design}
% ==========================================================

\subsection{Write the API first}

Write the code the user would write before writing the code that makes it work.
The calling code is the part that must be good, because it is the part users
see, learn, and depend on. Everything behind the interface can be rewritten
later; the interface itself cannot.

Designing outward from the implementation produces signatures that expose
internal structure rather than user intent. Designing inward from the call site
produces the opposite, and reveals awkwardness while it is still free to fix.

Judge a proposed interface by reading realistic usage rather than the
definition. Argument order, naming, and defaults that appear reasonable in a
signature often read badly in the script that calls them, and that script is
the true specification.

An interface that cannot be written cleanly is a signal that the underlying
concept is not yet clear. Resolve that before implementing, since implementation
effort tends to entrench whatever interface it was built to serve.



\subsection{Make the common case easy, and the uncommon case possible}
\label{sec:layers}

A useful design principle for scientific software is to provide the API in
\emph{layers of abstraction}. Each layer exposes the same underlying
capabilities but with different trade-offs between convenience and control.

Higher layers provide functions that solve common tasks with minimal user
effort. These functions often make reasonable assumptions, provide sensible
defaults, and combine multiple lower-level operations into a single call.
They are designed for the typical use case.

Lower layers expose smaller, more specialized primitives. These functions
require more explicit code from the user but provide greater flexibility,
composability, and control over the computational workflow.

Each layer should be built from the public primitives of the layer below, so
that the convenient functions are compositions of the controllable ones rather
than a second implementation of the same knowledge (Subsection~\ref{sec:dry}).

The relationship between abstraction layers, control, and convenience is
illustrated in Figure~\ref{fig:layered_api}. As abstraction layers increase,
convenience typically increases while direct control decreases.

\begin{figure}[h]
\centering
\begin{tikzpicture}[x=1cm,y=1cm]

\def\W{10}
\def\H{0.8}

\draw[->] (0,2.2) -- (\W,2.2);
\node[above] at (\W/2,2.2) {Layers};

\node[left] at (0,1.4) {Control};
\draw (0,1) rectangle (\W,1+\H);
\fill[black] (0,1) -- (0,1+\H) -- (\W,1) -- cycle;

\node[left] at (0,0.2) {Convenience};
\draw (0,-0.2) rectangle (\W,-0.2+\H);
\fill[black] (0,-0.2) -- (\W,-0.2+\H) -- (\W,-0.2) -- cycle;

\end{tikzpicture}
\caption{Layered API design. With increasing abstraction layers, direct
control decreases while convenience increases.}
\label{fig:layered_api}
\end{figure}



\subsection{Name axes, units, and conventions}
\label{sec:conventions}

Axis order, such as observations-by-variables versus
variables-by-observations or time-major versus space-major, the index base,
coordinate systems and reference frames, angle conventions, and physical units
are part of the meaning of every array. If they are not in the signature, the
documentation, or the returned object, users will infer them wrongly. The
package must not infer them from the data either
(Subsection~\ref{sec:no-guess}).

Do not rely on a comment or an example notebook to carry a convention. Name
the axes, state the units, and prefer the labeled data structures of the host
stack where they exist. Keep identifiers the caller supplied, such as sample
IDs, gene names, station codes, or node labels, intact and aligned with the
data through every operation.



\subsection{Numerics and defaults are part of the API}
\label{sec:numeric-defaults}

Changing a default normalization, initialization, regularization strength,
convergence tolerance, iteration limit, $\varepsilon$, reduction, or
tie-breaking rule can change published numbers. Such changes must be treated
as breaking changes of the API, not as internal tweaks.

Signatures are not the whole contract. Documented numerical defaults are part
of the scientific behavior of the package. If a default has to change, follow
the process for breaking changes in Subsection~\ref{sec:stable-api}.



\subsection{Design stable APIs}
\label{sec:stable-api}

API stability builds user trust, and in science it protects published
results. Avoid breaking changes that are not necessary. Keep the public API
small, since everything public is a commitment, and mark internal code as
private using the host stack's convention.

Number releases so that the version communicates compatibility, for example
with semantic versioning, and record every change of behavior in a changelog.
Changes of numerical behavior, including changes of defaults
(Subsection~\ref{sec:numeric-defaults}), must be listed explicitly in the
changelog, not summarized as fixes.

When a breaking change is necessary, deprecate before removing: emit a
deprecation warning (Subsection~\ref{sec:output}) that names the replacement,
keep the old behavior available for a stated period, and document the
migration path.



\subsection{Document the method, not the implementation}
\label{sec:document-method}

Documentation should describe what a function computes and how to use it, not
how the code is organized internally. Clear interfaces allow the
implementation to change without breaking user code.

In scientific software, however, the method is part of the behavior. The
documentation of a user-facing function should state the algorithm or
estimator, the approximations and assumptions it makes, its numerical
defaults, its policy for missing values (Subsection~\ref{sec:missing}), and
references to the literature that defines the method. Users need this
knowledge to interpret and report their results.

Implementation details, such as loop structure, internal data structures, and
private helper functions, do not belong in user documentation. The dividing
line is whether a change could alter a scientific result: if it could, it is
part of the method and must be documented; if it could not, it is
implementation.



% ==========================================================
\section{Explicit behavior and no side effects}
% ==========================================================

\subsection{Do what it says on the tin}
\label{sec:explicit}

A function should do what its name and arguments say, and nothing more.
Processing steps that change the data, such as filtering, normalization,
sorting, deduplication, or outlier removal, must not be hidden inside a
function named for a different operation: a user who calls a fitting
function should not discover later that the data were also filtered. Expose
such steps as named arguments, or provide them as separate functions the
user calls. Steps that are part of the definition of the method are
documented as such (Subsection~\ref{sec:document-method}).

This is the counterpart of Subsection~\ref{sec:no-guess}: a function must
neither read more into a call than it states nor do more than it asks.



\subsection{Don't guess}
\label{sec:no-guess}

The package must not resolve an ambiguity in the user's intent on the user's
behalf. A default is not a guess: it is fixed, documented, and independent of
the data. A guess applies heuristics to the data to decide what the call
means, such as transposing a matrix because of its shape, detecting whether
data are already normalized, inferring units or date formats from the values,
matching labels approximately, or repairing malformed input. If such an
interpretation could change the result, the user has to decide: it is a
scientific choice, and it belongs to the user. A guess also makes the meaning
of a call depend on the data, so the same code can do different things for
different datasets. When the guess is wrong, the result is wrong without
looking wrong (Subsection~\ref{sec:loud}).

Make the user state the choice through an argument, or raise an error that
lists the possible interpretations. Where automatic selection is genuinely
useful, offer it as an explicit option, such as a \texttt{method} argument
with the value \texttt{auto}, and record the choice in the returned result.
Heuristics may trigger warnings (Subsection~\ref{sec:output}), but must not
change behavior: announcing a guess does not make it the user's choice.


\subsection{Do not mutate caller data}
\label{sec:mutation}

A function must not overwrite, transpose, densify, drop missing values from,
or otherwise rewrite objects the caller still holds, unless the user opted
into an explicit in-place operation, such as a clearly named in-place function
or argument. Silent mutation makes it impossible to know which object still
holds the original data.

If a transformation is required, return a new object. The default is: leave
caller data untouched.



\subsection{Do not touch global state}
\label{sec:global-state}

Library code must not change process-wide settings. It must not set or reseed
the global random number generator, select hardware devices, change the number
of threads used by numerical libraries, change floating-point error handling,
change display options of data libraries, or rewrite plotting defaults. These
settings belong to the user. A package that calls \texttt{set.seed} in R or
\texttt{numpy.random.seed} in Python inside a function silently takes control
of the user's session.

A change that is strictly scoped to the package's own computation, and
restored before control returns to the caller, including when an error
occurs, may be made where the host stack provides a mechanism for it. How
the caller controls randomness instead is defined in
Subsection~\ref{sec:reproducibility}.



\subsection{The library follows the caller's data types}
\label{sec:follow-types}

If the caller passes data of a particular type, precision, or placement, keep
it. Library code must not silently densify sparse data, upcast single to
double precision, move data between the main processor and an accelerator, or
convert one array type of the host stack into another. The output should have
the same kind of type, precision, and placement as the input. Inputs that
carry no such choice, such as a plain list of numbers, may be converted to the
package's working type.

The caller's choices reflect deliberate trade-offs of memory, speed,
precision, and hardware. Following the incoming data preserves them; choosing
for the user breaks the workflows built on them.

Correctness outranks this rule (Subsection~\ref{sec:quality}). If the method
cannot produce a correct result in the incoming precision, raise an error that
names the required precision, rather than silently upcasting or silently
returning a degraded result (Subsection~\ref{sec:loud}); a user who wants the
conversion can make it explicitly. Using higher precision for small
intermediate quantities, such as an accumulator, is part of the method and
must be documented as such (Subsection~\ref{sec:document-method}); converting
the caller's data is not.
If the package cannot operate on a given type at all, it should say so with an
error rather than convert the data at hidden cost.



\subsection{No surprise cost}
\label{sec:cost}

Time, memory, and network use are part of the contract. A helper must not, as
a side effect, densify sparse data, materialize lazy data, download data, or
transfer a whole array to or from an accelerator.

Expensive steps should be explicit in the function the user calls, or
impossible to invoke by accident. A convenience wrapper that silently
materializes the full dataset is a bug, not a feature. Input validation has
costs as well; how they are balanced against correctness is defined in
Subsection~\ref{sec:validation-cost}.



\subsection{Importing the package has no side effects}
\label{sec:import}

Importing or loading the package must not download data, compile extensions,
print, select backends, configure logging, or read from or write to the
user's home directory. Users load libraries inside other tools, test suites,
and notebooks; work done at load time is hidden work, and it slows down and
destabilizes everything that depends on the package.

Put optional heavy setup behind an explicit function the caller invokes.
Loading the package should be cheap and silent.



\subsection{Raise, warn, or stay silent}
\label{sec:output}

Library code must not print. Standard output belongs to the user's scripts,
notebooks, and batch jobs. The package communicates through three channels,
and each situation calls for exactly one of them:

\begin{itemize}
  \item \textbf{Raise} when the call is invalid, or when continuing would
    return an untrustworthy result as if it were valid. The user has to act,
    and execution must not continue as if nothing happened.
  \item \textbf{Warn} when the call is valid but is probably not what the
    user intended, for example when a heuristic suggests that the data are
    not what the call assumes (Subsection~\ref{sec:no-guess}), or when the
    behavior of the call will change in a future version. Deprecations
    (Subsection~\ref{sec:stable-api}) belong here. Use the host stack's
    standard warning mechanism, so that users can filter warnings or turn
    them into errors.
  \item \textbf{Stay silent} when the call is valid and unremarkable.
\end{itemize}

A warning must not replace an error: if continuing would produce a wrong
result, raise (Subsection~\ref{sec:loud}). Warnings should name the argument
or condition that triggered them and say how to avoid them by changing the
call. Avoid emitting the same warning repeatedly, for example inside a loop.

Progress reports and diagnostic messages fall outside these three channels.
They should be opt-in, through an explicit argument or the host stack's
standard logging or messaging facility, so that the user decides whether and
where they appear.



% ==========================================================
\section{Input validation and errors}
% ==========================================================

\subsection{Garbage in, useful error out}
\label{sec:validation}

User-facing functions must validate their inputs at the point of entry:
types, shapes, bounds, units where they are declared, and the semantic
constraints of the method. Invalid inputs, unsupported combinations of
arguments, and dangerous settings must produce an error at the boundary, not
a numerical failure or a wrong result further downstream. Fail loudly and
early (Subsection~\ref{sec:loud}).

Error messages should be specific enough that users can recover without
searching through the full documentation: name the argument, state what was
expected and what was received, and where possible say how to fix the call.

Do not rely on documentation for correctness. Many users read documentation
only partially, skim examples, or rely on auto-completion and error messages
while using the package. Documentation explains correct usage; the code must
enforce the constraints that correctness depends on. Make misuse difficult and
failures obvious.

Implement each check once, in the validation module
(Subsection~\ref{sec:modules}), and call it from every entry point that needs
it (Subsection~\ref{sec:dry}). Private functions that are reached only through
a validated entry point should not repeat the checks.



\subsection{Keep validation proportionate}
\label{sec:validation-cost}

Validation has a cost, so the requirements for correctness and for cost
(Subsection~\ref{sec:cost}) can pull in different directions. Distinguish two
kinds of check:

\begin{description}
  \item[Metadata checks] inspect only types, shapes, precision, declared
    units, and combinations of arguments. They are cheap and must always run
    in user-facing functions.
  \item[Data checks] inspect the values themselves, for example for missing
    values, bounds, symmetry, or sortedness. They cost at least one pass over
    the data, and for data on an accelerator they can force the program to
    wait for the device.
\end{description}

Correctness outranks performance (Subsection~\ref{sec:quality}), so data
checks should run by default in user-facing functions, and their cost should
be documented. A user who has already validated the data may turn them off
through an explicit, clearly named argument. This follows the layered design
of Subsection~\ref{sec:layers}: the common case is safe, and the uncommon case
is possible.

Data checks must not themselves create a surprise cost
(Subsection~\ref{sec:cost}). For lazy or chunked data, express a check as part
of the same deferred computation, or apply it chunk by chunk.



\subsection{Missing values and NaN are explicit}
\label{sec:missing}

Each user-facing function must have a documented policy for missing values,
such as NaN, NA, or masked entries: reject, propagate, or mask. Silent
imputation, dropping, or filling is a scientific bug, not a convenience
(Subsection~\ref{sec:loud}): it replaces a missing measurement with a guess
(Subsection~\ref{sec:no-guess}).

If the input violates the policy, the function must raise an error that names
the missing values, rather than fail later with an unrelated numerical error.
Checking for missing values is a data check in the sense of
Subsection~\ref{sec:validation-cost}.



\subsection{Wear belt and braces}
\label{sec:assertions}

Tests check the package on the inputs someone thought of, but real data reach
paths that no test exercised. A bug that produces a silently wrong result is
the worst failure scientific software can have, so the key invariants of a
method should also be checked in the code itself, on every run. Such
assertions complement the tests; they do not replace them.

Assert what must hold if the code is correct and whose violation would
otherwise go unnoticed, such as a normalization, a conservation law, or the
alignment of identifiers after a reordering. These are often the same
invariants the tests check (Subsection~\ref{sec:invariants}). Keep assertions
cheap (Subsection~\ref{sec:cost}), and do not assert what holds by
construction or what validation already guarantees; mistakes of the caller
belong to validation (Subsection~\ref{sec:validation}).

A failed assertion means the package has a bug. It must raise, never warn,
and its message should say so and ask for a report, rather than tell the user
to change the call. Assertions must run in every configuration; where the
assertion mechanism of the host stack can be disabled, raise explicitly
instead.



% ==========================================================
\section{Numerical correctness}
% ==========================================================

\subsection{Prefer numerically stable formulations}
\label{sec:stability}

A formula that is correct in exact arithmetic can be badly wrong in floating
point. Choose formulations that remain accurate for the whole range of inputs
the package accepts, not only for well-behaved test data. Common patterns
include:

\begin{itemize}
  \item Avoid subtracting nearly equal quantities, which cancels significant
    digits. Rearrange the formula, or use dedicated functions such as
    \texttt{log1p}, \texttt{expm1}, and \texttt{hypot}.
  \item Work in log space for products of many small numbers, such as
    likelihoods, and compute $\log\sum_i e^{x_i}$ as
    $m + \log\sum_i e^{x_i - m}$ with $m = \max_i x_i$.
  \item Solve linear systems instead of forming explicit matrix inverses, and
    prefer factorizations suited to the structure of the problem.
  \item Accumulate long sums with care, for example by pairwise or
    compensated summation, or with a higher-precision accumulator
    (Subsection~\ref{sec:follow-types}).
\end{itemize}

Reuse the stable routines of mature libraries rather than reimplementing them
(Subsection~\ref{sec:core-domain}). Test behavior at the boundaries of the
input range, such as very small and very large magnitudes and nearly singular
inputs, as part of the invariant tests of Subsection~\ref{sec:invariants}.



\subsection{Report conditioning and convergence}
\label{sec:convergence}

Distinguish an unstable algorithm from an ill-conditioned problem. An unstable
algorithm is a defect of the package and must be fixed. An ill-conditioned
problem cannot be solved accurately by any algorithm; the package cannot fix
it, but it must not hide it. Where conditioning can be estimated at reasonable
cost, check it and report an ill-conditioned problem through the channels of
Subsection~\ref{sec:output}.

Iterative methods must not return a result that did not converge as if it had
(Subsection~\ref{sec:loud}). When the iteration limit is reached or the
convergence criterion fails, either raise, or return a result that carries an
explicit convergence status and warn. Each iterative function must document
which of the two it does, and its result should record diagnostics such as
the number of iterations and the final residual.



% ==========================================================
\section{Testing and reproducibility}
% ==========================================================

\subsection{Testing is not optional}
\label{sec:testing}

Scientific code is particularly prone to subtle errors that produce plausible
but wrong numbers. Every user-facing function must have tests, and every
release must pass the full test suite.

Tests should cover error paths as well as results: invalid inputs must be
shown to raise the documented errors (Subsection~\ref{sec:validation}). Every
fixed bug should gain a regression test that fails without the fix. Test the
range of input types the package claims to support, such as sparse,
single-precision, or lazy data (Subsection~\ref{sec:follow-types}); a code
path that is never exercised is not supported.

Run the test suite automatically on every change, across the versions of the
language and of the host-stack libraries that the package supports.



\subsection{Test invariants, not only examples}
\label{sec:invariants}

Passing on a few happy-path inputs is a weak scientific test. Prefer
invariants: conservation laws, symmetries, known analytic solutions and
limiting cases, consistency relations such as counts and normalizations, and
agreement with an independent implementation.

Examples are useful, but they do not prove that the method is correct. Encode
the scientific claims the software makes and test those claims directly.



\subsection{Size matters}
\label{sec:size-matters}

A method that works on test data but fails at the scale of real applications
solves no problem. Document the time and memory complexity of each
user-facing function in terms of the data size, and benchmark at the sizes of
the driving applications (Subsection~\ref{sec:applications}), not only on toy
inputs. Profile before optimizing, since bottlenecks are rarely where
intuition expects them, and track benchmarks to catch performance
regressions. Performance never outranks correctness or clarity
(Subsection~\ref{sec:quality}).


\subsection{Reproducibility: the same answer under stated conditions}
\label{sec:reproducibility}

Results must be reproducible. In these guidelines, reproducibility means
obtaining the same scientific answer, within a stated tolerance, from the same
inputs, parameters, and package version, on a documented software and
hardware stack. It does not mean identical bits forever: floating-point
results legitimately vary across linear algebra libraries, processors,
accelerators, and compiler settings.

Randomness must be controllable by the caller. Where the host stack supports
it, every function that uses random numbers should accept a random number
generator or seed argument that controls all randomness within that call.
Where the host stack's convention is a single global generator, draw from the
generator the caller controls, subject to Subsection~\ref{sec:global-state}.
The behavior when no generator or seed is supplied must be documented.

Prefer deterministic algorithms. Where nondeterminism is unavoidable, for
example because parallel reductions sum in varying order, document it and
ensure that results still agree within the stated tolerance.



\subsection{Compare floats with tolerance}
\label{sec:float-tests}

Tests should compare floating-point results within a tolerance, not bit for
bit. Bit-identical reference files treat platform noise as a logic bug and
make the test suite fragile.

Assert agreement within the tolerance that defines reproducibility
(Subsection~\ref{sec:reproducibility}). Use relative tolerances for quantities
of varying magnitude and absolute tolerances for quantities near zero.
Tolerances are constants in the sense of Subsection~\ref{sec:magic-numbers}:
derive each from the expected numerical error of the method, rather than
adjusting it until the tests pass on one machine.



\subsection{Flaky tests cry wolf}
\label{sec:flaky-tests}

A test that fails at random teaches people to ignore failures, and a test
suite whose failures are ignored checks nothing. Tests of stochastic methods
must therefore fix the random generator (Subsection~\ref{sec:reproducibility}),
so that every run gives the same outcome.

A fixed generator makes a test repeatable, not convincing: a test that passes
for only one seed can hide a bug. Where the claim is statistical, such as an
unbiased estimate or a coverage rate, assert it with bounds whose
false-failure rate is known and very small, for example a deviation of
several standard errors, so that the test would pass for other seeds as well.
The bound is a constant in the sense of Subsection~\ref{sec:magic-numbers}.
Never choose a seed because the test passes with it.



\subsection{Make every result traceable to a version}
\label{sec:traceability}

A result is reproducible only if the exact software that produced it can be
identified and obtained again, whether the result ends up in a publication, a
report, or a decision.

The package must expose its version at run time through the host stack's
standard mechanism, such as \texttt{\_\_version\_\_} in Python or
\texttt{packageVersion} in R. Every release should be tagged in version
control and archived so that the exact version can be retrieved later: in a
public archive with a persistent identifier, such as a DOI, for published
software, or in the organization's release archive for internal software.
Published software should also provide machine-readable citation metadata,
such as a \texttt{CITATION.cff} file, stating how to cite both the software
and the method.

Results should record the package version and the parameters that produced
them, including defaults that were not set explicitly: library code records
them in the result objects it returns, and the application layer records
them in the outputs and run metadata it writes. Changes of behavior between
versions are traced through the changelog (Subsection~\ref{sec:stable-api}).

\end{document}